\begin{lemma}
For every two-bite sample \(\omega\), every fixed set \(I\), and every base
vertex \(x\), the vertex \(x\) lies in one of the four cutoff classes:
\[
\noderef{TwoBiteSmallClass}(\omega,\varepsilon,I,x)
\quad\text{or}\quad
\noderef{TwoBiteMediumClass}(\omega,\varepsilon,I,x)
\quad\text{or}\quad
\noderef{TwoBiteLargeClass}(\omega,\varepsilon,I,x)
\quad\text{or}\quad
\noderef{TwoBiteHugeClass}(\omega,I,x).
\]
\end{lemma}
\begin{proof}
Let
\[
c=\bigl|\noderef{TwoBiteX}(\omega,I,x)\bigr|.
\]
Since \(\noderef{TwoBiteX}(\omega,I,x)\) is a filtered subset of \(I\), we
have \(0\le c\) and \(c\le |I|\).  Compare \(c\) successively with the huge,
large, and small cutoffs.  If
\(c>\noderef{TwoBiteHugeCutoff}(n)\), then \(x\) is huge.  Otherwise
\(c\le\noderef{TwoBiteHugeCutoff}(n)\); if
\(c>\noderef{TwoBiteLargeCutoff}(\varepsilon,n)\), then \(x\) is large.
Otherwise \(c\le\noderef{TwoBiteLargeCutoff}(\varepsilon,n)\); if
\(c>\noderef{TwoBiteSmallCutoff}(\varepsilon,n)\), then \(x\) is medium.
In the remaining case
\(0\le c\le\noderef{TwoBiteSmallCutoff}(\varepsilon,n)\), so \(x\) is small.
These are exactly the four defining alternatives.
\end{proof}
